\section{Does the gradient identify improving token replacements?}
The continuous score gradient gives a directional derivative at the current
embedding. A token replacement takes a finite step. We tested how often the
first-order approximation correctly identifies an improvement over that step.

For each ordered single-token replacement from $p$ to $q$, write $x(p)$ for
the implemented embedding and $F$ for the mean logistic score on the four
fixed training vectors. The predicted change and the measured change are
\[
 g = \nabla F(x(p))^\top [x(q)-x(p)], \qquad d = F(x(q))-F(x(p)).
\]
We classified a move as predicted improving when $g>10^{-12}$ and actually
improving when $d>10^{-12}$. For example, replacing the first token of
\texttt{(clear, clear)} with \texttt{json} gives $g=+0.006203$ and
$d=-0.020434$. The positive local slope overestimates the finite gain.

\begin{center}
\begin{tabular}{rrrrr}
\toprule
Length & True positive & False positive & False negative & True negative\\
\midrule
2 & 180 & 20 & 0 & 160\\
3 & 1620 & 133 & 0 & 1487\\
4 & 12960 & 851 & 0 & 12109\\
\bottomrule
\end{tabular}
\end{center}
Across the 29,520 directed replacements, the gradient flagged all 14,760
improvements and flagged 1,004 additional moves. Thus 14,760 of the 15,764
flagged moves improved the fixed training objective (93.63\%). These counts
cover the specified finite spaces exhaustively. They do not establish the same
property for longer prompts, other task vectors or language models.
A search algorithm using the sign as a filter would still need to score a
flagged move before accepting it. Whether this saves work depends on the
cost of obtaining gradients and on which prompts the search visits.

\paragraph{Reproduction.}
From the repository root run \texttt{python3 src/gradient\_audit.py}.
The script enumerates every position and every alternative token at each
starting prompt, and checks the total against $6^m(5m)$ for each length $m$.
